\section{More on metric spaces: connectedness}
A separation of a set is a partition into two nonempty relatively
open sets. A set is connected when it has no separation. A path is a
continuous map from an interval; path-connectedness asks for a path
between every two points of the set. Path-connected sets are connected,
but the converse fails. In this section $C(S)$ means the
\emph{bounded} continuous real-valued functions on $S\subseteq\R$,
and its distance is
\[
 d(f,g)=\|f-g\|_\infty=\sup_{x\in S}|f(x)-g(x)|.
\]
Here $\|h\|_\infty$ denotes the supremum of $|h|$. When $S$ is empty,
there is just one function and its distance from itself is taken to be
zero; the nontrivial arguments below concern nonempty $S$.

\exercise{22.1}{Two impossible continuous surjections}
\step{(a) A disconnected union of intervals.}
The interval $[0,1]$ is connected. By Theorem~22.2 a continuous image
must be connected. But $[0,1]\cup[2,3]$ has a separation into its two
nonempty pieces, obtained by intersecting with $(-\infty,3/2)$ and
$(3/2,\infty)$. Hence no continuous map from $[0,1]$ is onto that union.
\step{(b) A rational range.}
The interval $(0,1)$ is connected, whereas $\Q$ is not an interval:
for example, it contains $1,2$ but not $\sqrt2$ between them. Equivalently,
$\Q\cap(-\infty,\sqrt2)$ and $\Q\cap(\sqrt2,\infty)$ separate $\Q$.
\answer{Neither surjection exists, because a continuous image of a
connected set is connected.}

\exercise{22.2}{The unit circle is connected}
\goal Show the set $x_1^2+x_2^2=1$ is connected.
\step{Give a continuous parametrization with exactly that image.}
Define $\gamma:[0,2\pi]\to\R^2$ by
$\gamma(t)=(\cos t,\sin t)$. Both coordinates are continuous, so
Proposition~21.2 proves continuity of $\gamma$. The identity
$\cos^2t+\sin^2t=1$ puts its image on the circle. Conversely every
point of the unit circle has a polar angle $t\in[0,2\pi]$, and hence
is of this form.
\answer{The circle is the continuous image of the connected interval
$[0,2\pi]$, so Theorem~22.2 makes it connected.}

\exercise{22.3}{The closure of a connected set is connected}
\goal Prove connectedness of $\overline E$ from connectedness of $E$.
\step{Suppose there were a separation of the closure.}
Let open ambient sets $U,V$ give disjoint nonempty sets
$\overline E\cap U$, $\overline E\cap V$ whose union is $\overline E$.
Choose $u\in\overline E\cap U$. Since $U$ is an open neighborhood
of a closure point, $U\cap E\ne\varnothing$. The same argument gives
$V\cap E\ne\varnothing$.
\step{Restrict the separation.}
The sets $E\cap U$ and $E\cap V$ are disjoint, nonempty and relatively
open in $E$, and their union is $E$. They would separate $E$, contrary
to its connectedness.
\answer{No separation of $\overline E$ exists. For $E=\varnothing$,
the closure is empty and connected by definition.}

\exercise{22.4}{The closed topologist's sine curve}
\goal Study $E=\{(x,\sin(1/x)):0<x\le1\}$.
\step{(a) Find all new limit points.}
We claim
\[
 \overline E=E\cup(\{0\}\times[-1,1]).
\]
If $(x_n,\sin(1/x_n))\to(a,b)$, then $a\in[0,1]$ and $b\in[-1,1]$.
If $a>0$, continuity of $\sin(1/x)$ at $a$ forces $b=\sin(1/a)$,
so the limit is already in $E$. If $a=0$, it lies on the stated
vertical segment. Conversely, given $b\in[-1,1]$, choose an angle
$\theta$ with $\sin\theta=b$. For all sufficiently large integers $n$,
$x_n=1/(\theta+2\pi n)$ belongs to $(0,1]$ and tends to zero, while
$\sin(1/x_n)=b$. Thus $(0,b)\in\overline E$.
\step{(b) Establish connectedness before dealing with paths.}
The map $x\mapsto(x,\sin(1/x))$ is continuous on the connected
interval $(0,1]$. Its image $E$ is connected by Theorem~22.2, and
Exercise~22.3 makes its closure connected.
\step{(c) Explain why no path can leave the vertical segment.}
Suppose a path $\gamma(t)=(u(t),v(t))$, $0\le t\le1$, starts on the
vertical segment and ends in $E$. The coordinate functions are
continuous, $u(0)=0$, and $u(1)>0$. The zero set of $u$ is a nonempty
closed subset of $[0,1]$, so it has a largest element $t_0<1$.
We have $u(t_0)=0$ and $u(t)>0$ for $t>t_0$.

Choose a fixed $t_1>t_0$. Because $u(t_1)>0$, the intermediate value
theorem applied to $u$ on $[t_0,t_1]$ shows that it attains every
sufficiently small positive number. In particular, it attains the
numbers
\[
 a_n=\frac1{\pi/2+2\pi n},\qquad
 b_n=\frac1{3\pi/2+2\pi n}.
\]
Choose corresponding parameters $r_n,s_n\in(t_0,t_1]$. They tend
to $t_0$: for each $\eta>0$ with $t_0+\eta\le t_1$, the positive
continuous $u$ has a positive minimum on $[t_0+\eta,t_1]$; eventually
$a_n,b_n$ are smaller than that minimum. Thus their preimages cannot
lie in this latter interval.

For positive first coordinate, membership in $\overline E$ forces
$v(t)=\sin(1/u(t))$. Hence $v(r_n)=1$ and $v(s_n)=-1$. Both parameter
sequences approach $t_0$, contradicting continuity of $v$ there.
\answer{The closure is connected but is not path-connected. An
accumulating curve can touch a segment in its closure without supplying
any continuous path from that segment to the curve.}

\exercise{22.5}{Unions and intersections of connected sets}
\step{(a) A shared point joins two connected sets.}
Suppose $E\cap F\ne\varnothing$ but $E\cup F$ has a separation $A,B$.
Restricting the separation to $E$ shows $E$ lies entirely in one side,
since otherwise it would be disconnected. Similarly, $F$ lies entirely
in one side. A point of $E\cap F$ forces them to lie on the same side.
The other side would then be empty, a contradiction.
\step{(b) An intersection can be disconnected.}
Let
\[
 E=\{(x,y):x^2+y^2=1,\ y\ge0\},\qquad
 F=\{(x,y):x^2+y^2=1,\ y\le0\}.
\]
Each semicircle is a continuous image of a closed interval, hence
connected. But $E\cap F=\{(-1,0),(1,0)\}$ has two isolated points
and is disconnected: balls of radius $1/2$ about these points separate
it.
\answer{A nonempty intersection guarantees a connected union; it
does not guarantee a connected intersection.}

\exercise{22.6}{Why the supremum distance is a metric on bounded functions}
\step{(a) Verify finiteness and the metric axioms.}
If $|f|\le M$ and $|g|\le N$ on $S$, then
$0\le|f(x)-g(x)|\le M+N$. Thus $d(f,g)$ is finite and nonnegative.
It is symmetric because $|f(x)-g(x)|=|g(x)-f(x)|$.
If $d(f,g)=0$, then $|f(x)-g(x)|=0$ at every $x$, so $f=g$;
conversely identical functions have distance zero. Finally for every $x$,
\[
 |f(x)-h(x)|\le |f(x)-g(x)|+|g(x)-h(x)|
 \le d(f,g)+d(g,h).
\]
Taking the supremum over $x$ proves the triangle inequality.
\step{(b) Distinguish a finite metric from a convergence definition.}
Without boundedness, a pair such as $f(x)=x$, $g(x)=0$ on $\R$
would have $d(f,g)=\infty$, forbidden for an ordinary metric.
Uniform convergence itself only asks whether the errors $|f_n(x)-f(x)|$
are eventually small for all $x$ simultaneously. The functions can
be unbounded while their differences are uniformly small: $f_n(x)=x+1/n$
converges uniformly to $f(x)=x$ on $\R$.
\answer{Boundedness makes the supremum distance finite for every
pair in $C(S)$; it is not needed to state uniform convergence.}

\exercise{22.7}{A bounded sequence is a bounded continuous function on $\N$}
\goal Identify the metric space $B$ from Exercise~13.3 with $C(\N)$.
\step{Use discreteness to understand continuity.}
Under the usual distance inherited from $\R$, the ball of radius
$1/2$ about $n\in\N$ contains only $n$. Hence \emph{every} function
$f:\N\to\R$ is continuous: choose $\delta=1/2$ at each point, and
all tested values are identical to the center value.
\step{Match both the elements and their distances.}
The sequence $(x_1,x_2,\ldots)$ corresponds to the function $f(n)=x_n$.
A sequence is bounded exactly when its function is bounded. Moreover,
\[
 \sup_{n\in\N}|x_n-y_n|
 =\sup_{n\in\N}|f(n)-g(n)|.
\]
\answer{This correspondence is a bijection preserving distance, so
$B$ and $C(\N)$ are the same metric space after relabeling their elements.}

\exercise{22.8}{Evaluation at one point is uniformly continuous}
\goal For fixed $s_0\in S$, study $F(f)=f(s_0)$.
\step{Bound one error by the largest error.}
For $f,g\in C(S)$,
\[
 |F(f)-F(g)|=|f(s_0)-g(s_0)|\le\sup_{x\in S}|f(x)-g(x)|=d(f,g).
\]
Given $\varepsilon>0$, take $\delta=\varepsilon$. If $d(f,g)<\delta$,
the displayed inequality gives $|F(f)-F(g)|<\varepsilon$.
\answer{Evaluation is $1$-Lipschitz and therefore uniformly continuous.}

\exercise{22.9}{A straight line in a function space}
\goal Prove uniform continuity of $F(t)=tf+(1-t)g$ as a map
$\R\to C(S)$.
\step{First check that the values belong to the claimed space.}
For fixed $t$, the function $tf+(1-t)g$ is continuous and is bounded
by $|t|\|f\|_\infty+|1-t|\|g\|_\infty$. Thus $F(t)\in C(S)$.
\step{Compute the distance between two values.}
The difference is $F(t)-F(s)=(t-s)(f-g)$, so
\[
 d(F(t),F(s))=|t-s|\,\|f-g\|_\infty.
\]
Let $M=\|f-g\|_\infty$. If $M=0$, $F$ is constant. Otherwise
$\delta=\varepsilon/M$ makes $|t-s|<\delta$ imply
$d(F(t),F(s))<\varepsilon$.
\answer{The map is Lipschitz with constant $\|f-g\|_\infty$.
The estimate works on the whole real line.}

\exercise{22.10}{Translation is continuous in the supremum metric}
\goal For bounded uniformly continuous $f:\R\to\R$, prove uniform
continuity of $F(x)=f_x$, where $f_x(y)=f(x+y)$.
\step{Check each translated function.}
Translation preserves continuity, and $|f_x(y)|\le\|f\|_\infty$.
Thus $f_x\in C(\R)$ for every real $x$.
\step{Keep the error bound uniform in the variable being translated.}
Fix $\varepsilon>0$. Uniform continuity of $f$ supplies $\delta>0$
such that $|u-v|<\delta$ implies
$|f(u)-f(v)|<\varepsilon/2$ for all real $u,v$. If $|x-z|<\delta$,
then for \emph{every} $y$,
\[
 |f_x(y)-f_z(y)|=|f(x+y)-f(z+y)|<\varepsilon/2,
\]
because $|(x+y)-(z+y)|=|x-z|$. Taking a supremum yields
$d(F(x),F(z))\le\varepsilon/2<\varepsilon$.
\answer{The same $\delta$ works for all $x,z$, proving uniform continuity.
Using $\varepsilon/2$ avoids losing the strict inequality at a supremum.}

\exercise{22.11}{The closed unit ball and connectedness in $C(S)$}
\goal Let $\mathcal E=\{f\in C(S):\|f\|_\infty\le1\}$.
\step{(a) Its complement is open.}
Suppose $f\notin\mathcal E$. Then some $x_0\in S$ has $|f(x_0)|>1$;
otherwise $1$ would be an upper bound for all $|f(x)|$. Put
$r=(|f(x_0)|-1)/2>0$. If $d(f,g)<r$, then
\[
 |g(x_0)|\ge|f(x_0)|-|g(x_0)-f(x_0)|
 >|f(x_0)|-r>1.
\]
Thus every $g$ in this metric ball is also outside $\mathcal E$.
The complement is open, so $\mathcal E$ is closed.
\step{(b) The entire space is path-connected.}
For any $f,g\in C(S)$, Exercise~22.9 gives a continuous path
$t\mapsto(1-t)f+tg$, $0\le t\le1$, from $f$ to $g$. Thus $C(S)$
is path-connected, and Theorem~22.5 makes it connected.
\step{(c) The same paths stay in the unit ball.}
If $f,g\in\mathcal E$ and $0\le t\le1$, then at every $x$,
\[
 |(1-t)f(x)+tg(x)|\le(1-t)|f(x)|+t|g(x)|\le1.
\]
The path therefore lies in $\mathcal E$.
\answer{$\mathcal E$ is closed and path-connected, hence connected;
$C(S)$ is connected as well.}

\exercise{22.12}{A topology determined by finitely many evaluations}
\goal Verify the topology axioms for the specified pointwise neighborhoods.
Write
\[
 N(f;F,\varepsilon)=\{g\in C(S):
               |g(x)-f(x)|<\varepsilon\text{ for all }x\in F\},
\]
where $F\subseteq S$ is finite and $\varepsilon>0$.
\step{(a) Recall what must be proved.}
Discussion~13.7 describes the characteristic properties of open sets:
the empty set and the entire space are open, arbitrary unions of open
sets are open, and finite intersections of open sets are open. These
properties define a topology even when no metric has been specified.
Here a set $U$ is declared open if each $f\in U$ has some
$N(f;F,\varepsilon)\subseteq U$.
\step{(b) Check all three properties.}
The empty set is open because it has no points for which the requirement
could fail. The whole space is open since every such neighborhood lies
in it (one can take $F=\varnothing$). If $f\in\bigcup_\alpha U_\alpha$,
then $f$ belongs to some open $U_{\alpha_0}$, which supplies a
neighborhood contained in $U_{\alpha_0}$ and hence in the union.

For a finite intersection $U_1\cap\cdots\cap U_m$ and a point $f$ in
it, choose $N(f;F_j,\varepsilon_j)\subseteq U_j$ for every $j$.
The set $F=F_1\cup\cdots\cup F_m$ is finite, and
$\varepsilon=\min_j\varepsilon_j>0$. A function in
$N(f;F,\varepsilon)$ satisfies every constraint for every $U_j$,
so this neighborhood lies in the intersection. The empty finite
intersection is the whole space, already covered.
\step{Why these are pointwise rather than uniform neighborhoods.}
Each neighborhood restricts only finitely many values of a function.
A small supremum-metric ball restricts values at \emph{every} point.
For example, on $S=\R$ any finite set of prescribed points misses
some $x_0$; a continuous bounded bump supported near $x_0$ can be zero
at every prescribed point yet have supremum $2$. Thus no finite-point
neighborhood of zero is contained in the supremum ball of radius $1$.
\answer{The declared sets satisfy every topology axiom. The text's
additional blanket remark about nonmetrizability needs qualification:
for finite $S$ this is the usual finite-dimensional topology, and for
countable $S=\{s_1,s_2,\ldots\}$ it is metrized by
$\sum_{j\ge1}2^{-j}\min(1,|f(s_j)-g(s_j)|)$. The requested topology
verification is valid for every $S$.}

\exercise{22.13}{A connected graph with an oscillatory endpoint}
\goal Analyze the graph with the oscillating branch at positive $x$
and the added origin.
\step{Clarify the printed notation.}
The exercise prints $f(x)=(x,\sin(1/x))$ for $x>0$ and $f(0)=0$.
Read literally, this is a map to $\R^2$ with $f(0)=(0,0)$, and its
graph lies in $\R^3$. The frequently intended scalar formulation is
$f(x)=\sin(1/x)$, $f(0)=0$. Both give the same connectedness answer:
the literal graph is homeomorphic to
\[
 H=\{(x,\sin(1/x)):x>0\}\cup\{(0,0)\}
\]
by the map $(x,y)\mapsto(x,x,y)$, whose inverse on its image is
$(a,b,c)\mapsto(a,c)$. Both maps are continuous coordinate maps.
Thus we prove the assertions for $H$.
\step{The added origin does not disconnect the branch.}
The branch $A=\{(x,\sin(1/x)):x>0\}$ is a continuous image of the
interval $(0,\infty)$, hence connected. Also $(0,0)\in\overline A$,
since $(1/(n\pi),0)\to(0,0)$. If $H$ had a separation, connectedness
would put all of $A$ on one side. The other side would consist of the
origin and be relatively open in $H$, which would provide a
neighborhood of the origin missing $A$. This contradicts its being
a closure point. Hence $H$ is connected.
\step{It has no path from the origin to the branch.}
The proof of Exercise~22.4(c) applies without change. A proposed path
has a last parameter $t_0$ at which its first coordinate is zero.
Immediately after $t_0$ its first coordinate is positive and tends
to zero. It must attain the values $a_n$ and $b_n$ from that exercise
at parameters tending to $t_0$, forcing its second coordinate to take
both $1$ and $-1$ arbitrarily close to $t_0$. That contradicts
continuity. The restriction $x\le1$ was not used in that argument.
\step{Continuity of the original function fails at zero.}
Along $x_n=1/(\pi/2+2\pi n)$, the sine coordinate is $1$, not tending
to its value $0$ at zero. The sequential criterion therefore rules out
continuity at zero in either formulation; continuity holds for $x>0$.
\answer{The graph is connected but not path-connected, and the
function is not continuous on $[0,\infty)$.}

\exercise{22.14}{Small successive changes force an interval of cluster values}
\goal For bounded $(s_n)$ with $s_{n+1}-s_n\to0$, show its set $S$
of subsequential limits is a connected interval.
\idea If the sequence repeatedly visits both sides of a level $c$,
it must cross $c$ repeatedly. Small jumps make some crossing terms
arbitrarily close to $c$.
\step{Fix any level between two subsequential limits.}
Take $a,b\in S$ with $a<c<b$. For each integer $j\ge1$, choose $N_j$
large enough that $|s_{n+1}-s_n|<1/j$ for all $n\ge N_j$. Since $a$
and $b$ are subsequential limits, there are arbitrarily late terms
below $c$ and arbitrarily late terms above $c$. Indeed, use errors
smaller than $(c-a)/2$ and $(b-c)/2$ in their defining subsequences.
\step{Locate crossings farther and farther out.}
Choose $p_j>\max(N_j,q_{j-1})$ with $s_{p_j}<c$, and then choose
$q_j>p_j$ with $s_{q_j}>c$ (start with $q_0=0$). Let $k_j$ be the
first index after $p_j$ for which $s_{k_j}\ge c$; it exists because
$q_j$ is one such index. Then
\[
 s_{k_j-1}<c\le s_{k_j},\qquad
 0\le s_{k_j}-c<s_{k_j}-s_{k_j-1}<1/j.
\]
The chosen indices strictly increase, so $(s_{k_j})$ is a subsequence
and converges to $c$. Thus $c\in S$.
\step{Conclude connectedness.}
The preceding argument shows $S$ contains every number strictly
between any two of its points, which is exactly the interval property.
Boundedness and Bolzano--Weierstrass ensure $S$ is nonempty; the
cluster-set theorem also makes it closed and bounded. Thus it is a
closed bounded interval, possibly a singleton, and intervals are connected.
\answer{The cluster values cannot have a gap because vanishing step
sizes prevent repeated crossings from jumping across such a gap.}
